%% ==============================================================================
%% content/04-resultats.tex
%% Chapitre 4 : Résultats expérimentaux et validation
%% ==============================================================================
\chapter{Résultats Expérimentaux et Validation}
\label{chap:resultats}

\section{Protocole Expérimental et Métriques}
\label{sec:resultats-protocole}

Pour valider l'efficience et la robustesse de notre architecture,
nous avons conduit une campagne de bancs d'essai sur plusieurs jeux de données représentatifs.
Les grandeurs physiques et les métriques de performance sont systématiquement formulées
à l'aide du paquet \code{siunitx}, assurant le respect des normes métrologiques internationales.
À titre d'exemple, la puissance crête mesurée est de \qty{142.8}{\kilo\watt},
tandis que la vitesse de convergence moyenne s'établit à \qty{1.45e-3}{\second\per\iter}.

\begin{definitionbox}[Indice d'efficience typographique ($\eta_T$)]
\label{def:efficience}
Soit $N_p$ le nombre total de pages d'un document et $T_c$ le temps de compilation complet avec Biber.
L'indice d'efficience typographique est défini par le ratio :
\begin{equation}
  \eta_T = \frac{N_p}{T_c} \quad \left[\text{pages}\cdot\unit{\second^{-1}}\right]
\end{equation}
Un ratio élevé témoigne d'une compilation optimisée sans goulets d'étranglement dans le préambule.
\end{definitionbox}

\section{Analyse Comparative des Performances}
\label{sec:resultats-analyse}

Les données quantitatives obtenues lors des campagnes de test
sont récapitulées dans le \cref{tab:resultats-synthese}.
Comme préconisé par les bonnes pratiques typographiques,
ce tableau utilise exclusivement les filets horizontaux du paquet \code{booktabs}
(\code{\textbackslash toprule}, \code{\textbackslash midrule}, \code{\textbackslash bottomrule})
et proscrit tout filet vertical afin de favoriser le balayage visuel continu de la donnée.

\begin{table}[htbp]
  \centering
  \caption{Synthèse comparative des performances selon la configuration de compilation}
  \label{tab:resultats-synthese}
  \begin{tabularx}{\textwidth}{@{} l r r r X @{}}
    \toprule
    \textbf{Configuration} & \textbf{Mémoire (\unit{\mega\byte})} & \textbf{Temps (\unit{\second})} & \textbf{Score $\eta_T$} & \textbf{Observations} \\
    \midrule
    Modèle Monolithique   & \num{248.5} & \num{8.42} & \num{12.4} & Avertissements de fontes et collision de paquets \\
    Extensions disparates & \num{310.2} & \num{9.15} & \num{11.5} & Redondance des inclusions de macros \\
    \textbf{Notre classe (Écran)}   & \num{182.4} & \num{4.18} & \num{25.1} & Navigation fluide, hyperliens dynamiques \\
    \textbf{Notre classe (Papier)}  & \num{185.0} & \num{4.22} & \num{24.9} & Géométrie BCOR compensée, contraste maximal \\
    \bottomrule
  \end{tabularx}
\end{table}

Ces résultats corroborent les observations de \textcite{dupont2024methodologie} :
la rationalisation des couches d'importation permet de diviser par deux le temps de compilation
tout en garantissant une totale immunité face aux régressions visuelles.
Comme le démontrent également les travaux de \textcite{goodfellow2016deep} et \textcite{vaswani2017attention},
la rigueur dans la structuration des composants logiciels et des pipelines de traitement
constitue un facteur déterminant pour la reproductibilité expérimentale.

\section{Discussion et Synthèse}
\label{sec:resultats-discussion}

L'évaluation qualitative met en lumière trois bénéfices majeurs pour l'utilisateur :
\begin{itemize}
  \item \textbf{Lisibilité du code source :} un nouvel intervenant identifie instantanément
  où ajouter un chapitre, une macro ou une entrée bibliographique sans risque d'effets de bord.
  \item \textbf{Polyvalence du livrable :} générer la version finale prête pour l'imprimeur
  ou la version numérique pour le jury ne nécessite qu'un seul basculement d'option.
  \item \textbf{Stabilité pérenne :} l'absence d'extensions obsolètes
  garantit une compatibilité totale avec les distributions \TeX{} Live contemporaines et futures.
\end{itemize}

